%% Преамбула TeX-файла

% 1. Стиль и язык
\documentclass[utf8, 12pt]{G7-32} % Стиль (по умолчанию будет 14pt)
%Работа с картинкой - нужно для герба - Krage56
\usepackage{graphicx}
\usepackage{float}
\usepackage{wrapfig}

%Работа с таблицами - Krage56
\usepackage{ragged2e}
\usepackage{booktabs}
\usepackage{amsmath}

% Остальные стандартные настройки убраны в preamble.inc.tex.
\sloppy

% Настройки стиля ГОСТ 7-32
% Для начала определяем, хотим мы или нет, чтобы рисунки и таблицы нумеровались в пределах раздела, или нам нужна сквозная нумерация.
\EqInChapter % формулы будут нумероваться в пределах раздела
\TableInChapter % таблицы будут нумероваться в пределах раздела
\PicInChapter % рисунки будут нумероваться в пределах раздела

% Добавляем гипертекстовое оглавление в PDF
\usepackage[
bookmarks=true, colorlinks=true, unicode=true,
urlcolor=black,linkcolor=black, anchorcolor=black,
citecolor=black, menucolor=black, filecolor=black,
]{hyperref}

% Изменение начертания шрифта --- после чего выглядит таймсоподобно.
% apt-get install scalable-cyrfonts-tex

\IfFileExists{cyrtimes.sty}
    {
        \usepackage{cyrtimespatched}
    }
    {
        % А если Times нету, то будет CM...
    }

\usepackage{graphicx}   % Пакет для включения рисунков
\DeclareGraphicsExtensions{.jpg,.pdf,.png}
% С такими оно полями оно работает по-умолчанию:
% \RequirePackage[left=20mm,right=10mm,top=20mm,bottom=20mm,headsep=0pt]{geometry}
% Если вас тошнит от поля в 10мм --- увеличивайте до 20-ти, ну и про переплёт не забывайте:
\geometry{right=20mm}
\geometry{left=30mm}



% Произвольная нумерация списков.
\usepackage{enumerate}

\setcounter{tocdepth}{1} %Подробность оглавления
%4 это chapter, section, subsection, subsubsection и paragraph
%3 это chapter, section, subsection и subsubsection
%2 это chapter, section, и subsection
%1 это chapter и section
%0 это chapter.

\begin{document}

\frontmatter % выключает нумерацию ВСЕГО; здесь начинаются ненумерованные главы: реферат, введение, глоссарий, сокращения и прочее.
\begin{center} 
    \begin{figure}[h]
        \centering
        \includegraphics[width=30mm, height=30mm]{inc/img/Gerb.png}
        \label{fig:Gerb}
        \begin{tabular}{p{15cm}}
            \begin{center}
                \small{МИНОБРНАУКИ РОССИИ}\\
                \small{Федеральное государственное бюджетное образовательное учреждение высшего образования}\\
                \textbf{«МИРЭА – Российский технологический университет»}\\
                \textbf{РТУ МИРЭА}\\
                \rule{\linewidth}{1pt}\\[-0.25cm]
                \small{Институт кибернетики}\\
                \small{Кафедра }\\[-0.45cm]
                \rule{\linewidth}{.1pt}
            \end{center}
        \end{tabular}
    \end{figure}
    
\Large РЕФЕРАТ \\[4pt] % название работы, затем отступ 0,6см
\Largeпо «Название предмета»\\[4pt]
\Largeна тему: «Название темы»
\end{center}
\vspace{8pt}
\large{\begin{center}
    \begin{tabular}{p{4.99cm}p{4cm}p{.01cm}p{4cm}}
    Обучающийся & &\multicolumn{2}{c}{Данные обучающегося}\\ 
        \cline{2-2}\cline{4-4}
    & \multicolumn{1}{c}{\small{\textit{Подпись}}} & & \multicolumn{1}{c}{\small{\textit{Фамилия Имя Отчество}}} \\
    Шифр & \multicolumn{3}{l}{Шифр студ. билета}\\\cline{2-2}
     & & & \\
    Группа & \multicolumn{3}{l}{К***-**-**}\\\cline{2-2}
     & & & \\
    Руководитель работы & &\multicolumn{2}{c}{Данные руководителя}\\ 
        \cline{2-2}\cline{4-4}
    & \multicolumn{1}{c}{\small{\textit{Подпись}}} & & \multicolumn{1}{c}{\small{\textit{Фамилия Имя Отчество}}} \\
    \end{tabular}
\end{center}}
\vfill
\begin{center} 
\large Москва 2020
\end{center} 

\thispagestyle{empty}

\thispagestyle{empty}
\setcounter{page}{0}
\tableofcontents
\clearpage


%\Introduction

Текст вступления

\mainmatter

\chapter{Первый раздел}
\label{cha:ch_1}


\chapter{Второй раздел}

\chapter{Третий раздел}
\label{cha:ch_3}

\backmatter %% Здесь заканчивается нумерованная часть документа и начинаются ссылки и
            %% заключение

\Conclusion

\include{5-biblio}

\appendix   % Тут идут приложения

\include{61-appendix}

\end{document}
